\section{A positive example beyond every one-sided jet extension}

\subsection{O-025: strictness of the two-sided order-ideal criterion}
Set
\[
 m=2^{200},\qquad r=m^2,\qquad n=m^{25},\qquad d=3n,
             \qquad B_0=1000000,
\]
and define
\[
 \lambda=(nd-26r-1999928,\ B_0+6,\ B_0,\ 26^{r-3}).
                                                               \tag{O-025.1}
\]
Then
\[
                  4\leq q^m_\lambda(d[n])
                                \leq K_n(\lambda)=a_\lambda(d[n]).
                                                               \tag{O-025.2}
\]
The two-sided construction M-019 proves the ambient
determinant equality. No choice of base jet order and
finite one-sided order ideal in M-018 can prove that
equality at size $n$ for this support. This compares
all of M-018's admissible choices, not just a particular
canonical ideal or a prefix simplex.
The exact values of the multiplicities remain unspecified.

\paragraph{The two-sided extension fits.}
Use the canonical shaving O-024 with $b=26$.
The shaved support has only two nonzero bounds:
\[
                       (B_0-20,\ B_0-26).
\]
It is contained in the sum of the two axis ideals
\[
 \mathcal A=\{1,x_2,\ldots,x_2^{B_0-20}\},\qquad
 \mathcal B=\{1,x_3,\ldots,x_3^{B_0-26}\}.
\]
Therefore $\Gamma_\lambda\subseteq\mathcal A+\mathcal B+D_{26}$.
At this perfect-square length, both
$\lceil\sqrt r\rceil$ and $\lceil\sqrt{r-1}\rceil$ equal $m$.
With
$T=m\binom{r+11}{12}$, the two-sided size is
\[
          N=(|\mathcal A|+|\mathcal B|)T
                =(2B_0-44)T=1999956T<n.                \tag{O-025.3}
\]
The final inequality follows from $1999956<C_*=478126755$
and M-017's normalized estimate. M-019 gives $K_n=a_\lambda$.
No uniform dominance relation between the two criteria
at other, nonsquare ambient lengths is needed here.

\paragraph{Four positive permanent copies.}
Multiply the four degree-six binary classes of M-012 by
\[
                      W_r^{13}\,W_3^{499987}\,
                             c_0^{\,d-13r-1499967}.     \tag{O-025.4}
\]
The last exponent is nonnegative, and the Hessian argument
O-011 makes all factors simultaneously nonzero on the
permanent closure. Its coordinate algebra is an integral
domain, so all four classes remain independent.
The total degree is
$6+13r+3\cdot499987+(d-13r-1499967)=d$.
The second weight is
$6+26+2\cdot499987=1000006$, the third is
$26+2\cdot499987=1000000$, and the later weights are $26$.
Their first weight is the remaining part in (O-025.1).
This proves $q^m_\lambda\geq4$ and hence (O-025.2).

\paragraph{Excluding every one-sided choice.}
Suppose, for a contradiction, that M-018 applies with
some even base order $b=2\ell+2$ and a nonempty finite
order ideal $A$:
\[
 \Gamma_\lambda\subseteq A+D_b,\qquad
             T_b(1+|A|)\leq n,\qquad
             T_b=m\binom{r-1+\ell}{\ell}.
                                                               \tag{O-025.5}
\]
Taking more ambient variables can only increase this
core size. Monomials using those additional variables
cannot participate in a nonnegative decomposition of
an exponent supported in the first $r$ variables, so
discarding them from $A$ retains every decomposition
needed below. Thus (O-025.5) is the weakest relevant
size bound even if a larger ambient length was chosen.

If $b\geq28$, then $\ell\geq13$ and
\[
 T_b(1+|A|)\geq2m\binom{r+12}{13}
                       \geq\frac{2m^{27}}{13!}>n.
                                                               \tag{O-025.6}
\]
This already contradicts feasibility.

If $b=26$, then $T_b\geq n/12!$, so
$|A|<12!=479001600$. There are $r-3>12!$ suffix
variables $x_4,\ldots,x_r$. At least one of their unit
monomials, say $x_j$, is missing from $A$.
Downward closure implies that no member of $A$ uses $x_j$.
For every two-variable monomial $\nu(x_2,x_3)$ of total
degree at most $B_0-26=999974$, the monomial
$x_j^{26}\nu$ lies in $\Gamma_\lambda$.
In a decomposition of its exponent as $\alpha+\eta$,
with $\alpha\in A$ and $|\eta|\leq26$,
the entire $x_j^{26}$ must belong to $\eta$.
It exhausts that degree budget, so $\eta=26e_j$ and
$\alpha=\nu$. Consequently $A$ contains the entire
two-variable simplex of degree $999974$, and
\[
                    |A|\geq\binom{999976}{2}
                        =499975500300>12!,
                                                               \tag{O-025.7}
\]
a contradiction.

It remains to treat $2\leq b\leq24$. Then
$0\leq\ell\leq11$ and
$k=26-b=24-2\ell\geq2$.
Put $R'=r-3$. The support $\Gamma_\lambda$ contains
every squarefree degree-$26$ monomial in those last
$R'$ variables. In its decomposition $\gamma=\alpha+\eta$,
the exponent $\alpha$ is squarefree and has degree at
least $k$. By downward closure, $A$ contains a
squarefree $k$-submonomial of $\gamma$.
Thus the squarefree degree-$k$ members of $A$ meet every
$26$-subset of the suffix variables. Double-counting
the incidences gives
\[
 |A|\geq
 \frac{\binom{R'}{26}}{\binom{R'-k}{26-k}}
       =\frac{\binom{R'}k}{\binom{26}k}
       \geq\frac{r^k}{2^k\,26!}.                       \tag{O-025.8}
\]
The final inequality follows because every numerator
factor in $\binom{r-3}k$ is at least $r/2$;
the additional factor $(26-k)!$ was harmlessly omitted.
Also
$T_b\geq m r^\ell/\ell!$ and $n=m r^{12}$.
It follows that
\[
 \frac{T_b|A|}{n}
       \geq\frac{r^{12-\ell}}{2^k\ell!\,26!}>1.
                                                               \tag{O-025.9}
\]
Indeed $12-\ell\geq1$ and
$2^k\ell!\,26!\leq2^{24}11!\,26!<2^{139}<r=2^{400}$.
This contradicts (O-025.5) for every remaining base order.
All cases are exhausted, proving the asserted strict
separation from every one-sided M-018 choice.

\paragraph{Other established failures remain.}
The full support (O-022.6) is contained in the present
support, and every present row cap is at least its
O-022 counterpart. Thus O-022's exact generic
degree-$26$ complexity argument still excludes every
F-018 exact support cover, while its separate dimension
arguments exclude every entire one-cut family and the
whole point-jet target.
The coordinate-box size failure also persists.
The shape meets the BIP staircase and A-025's joint gate.
The shape-specific Chow threshold is at least
$d-1=3n-1>n$, and the uniform Chow and Grenet thresholds
fail as before. Its length is $m^2$, its tail is
$26r+1999928\leq md$, and its third row is larger than two.
It remains outside every one-sided jet exclusion by the
complete argument (O-025.5)--(O-025.9).
The new two-sided comparison therefore reduces the
actual remaining gate on a shape with at least four
permanent copies. It does not produce a multiplicity
obstruction; (O-025.2) rules one out here.

\paragraph{Geometric dictionary and scope.}
The two axis ideals are finite free integral torus-weight
modules under SGA~3 I, Proposition~4.7.3, pp.~35--36.
The constant bridge and the two attached families in
M-019 give their actual universally injective coordinate
pullback. F-016 supplies the full Schur-copy support
through SGA~6 VI, Theorem~3.3, pp.~392--393, and
Macdonald I(4.3), p.~63, Appendix~A(5.3)--(5.6),
pp.~155--156. The counting argument (O-025.8)
is a positive incidence count in the suffix monomial
basis, not an alternating virtual identity.
The exact representation barrier inherited from O-022
and its border dimension barriers retain their distinct
scopes. No general saturation or unresolved evaluation-rank
comparison is inferred.
